% paper/sections/introduction.tex
% Target: 0.9–1.1 pages (sigconf two-column ≈ 700–900 words body text)
% Style: WSDM framing — network economics first, learning as instrument.
%        Structure: P1 viral marketing+Babaei | P2 deep-IM gap (1 sentence) |
%                   P3 our approach | P4 budget realism | P5 C1–C4 bullets | P6 roadmap
% Every number has a %SOURCE comment. Phrases follow style_notes.md.
%
% ANONYMIZATION: cite WSDM 2027 prior work as \cite{anon_wsdm27}.
%
% NUMBER SOURCES:
%   462.6 / 460.0   % SOURCE: CLAUDE.md; results/logs/comparison_20260623_131609.json
%   p<0.0001 20/20  % SOURCE: CLAUDE.md §Session State
%   +8.6% n=2000    % SOURCE: results/logs/paper_gen_updated.json (FF n=2000)
%   +40.5% Rice-FB  % SOURCE: results/logs/paper_gen_updated.json (Rice-FB n=443)
%   99.8 / 23.6 k=3 % SOURCE: dp_v3_full_curve_merged.json / dp_upgrade_eval.json
%   305.6 / 7.4 k=1 % SOURCE: unified_sweep.json / dp_upgrade_eval.json

\section{Introduction}
\label{sec:intro}

% ── P1: Viral marketing + discounting + Babaei ────────────────────────────────
A firm selling a product through a social network faces a classical tension:
deep discounts convert hesitant buyers and trigger peer-to-peer
diffusion,  but they erode margin on every unit sold.
\citet{babaei2013revenue} formalized this trade-off as a
\emph{revenue maximization via discounting} problem: the seller sequentially
targets a set of buyers and offers each a personalized price, exploiting the
fact that a buyer's willingness-to-pay rises as more of her social contacts
have already purchased.
Their framework showed that calibrated discounts strictly dominate the
free-seeding paradigm of classical influence maximization~\citep{kempe2003maximizing}
in terms of total revenue.
The practical limitation is that they rely on two independent, hand-crafted
steps — greedy seed selection by degree centrality, then discount assignment via a
simple degree rule — designed separately and applied in sequence.
The coupling between \emph{which} buyer to target next and
\emph{what price} to offer is never exploited.

% ── P2: Deep IM gap (one sentence) ───────────────────────────────────────────
Deep reinforcement learning for influence maximization — notably
S2V-DQN~\citep{dai2017learning} and ToupleGDD~\citep{touplegdd} — achieves
strong results on the seed-selection subproblem but optimizes
\emph{diffusion reach}, not revenue, and assumes zero production cost and no pricing
decisions.

% ── P3: Our approach ─────────────────────────────────────────────────────────
Here, we close both gaps.
We extend the graph RL framework of our prior work~\citep{anon_wsdm27} to the
joint seed-and-pricing problem.
At each decision step a two-layer GraphSAGE encoder produces node embeddings
that capture both structural centrality and the current influence state of
every buyer.
An LSTM then compresses the full episode history — past offers, accept/reject
outcomes, running revenue — into a context vector that is concatenated with
the per-node embedding before scoring.
A~joint policy head then selects the next buyer (discrete) and outputs a
personalized discount $d_t \in [0,1]$ (continuous), trained first by imitation
of the \citet{babaei2013revenue} greedy expert and then refined by REINFORCE.
Trained only on Forest-Fire graphs with $n \le 440$ nodes, the model achieves
$462.6$ %SOURCE: CLAUDE.md; comparison_20260623_131609.json
revenue on the standard $n{=}1000$ test graph versus $460.0$ for the best
hand-crafted baseline ($p{<}0.0001$, %SOURCE: CLAUDE.md §Session State
$20/20$ seeds).
Zero-shot generalization: $+8.6\%$ at $n{=}2000$ %SOURCE: paper_gen_updated.json
and $+40.5\%$ on Rice-Facebook. %SOURCE: paper_gen_updated.json

% ── P4: Realism gap → budget formulation ─────────────────────────────────────
A second limitation of all prior work — including the unconstrained model
above — is the assumption of infinite capital: the firm can approach any buyer
at any cost.
In practice, a seller starts with finite budget $B_0 = k \cdot c$ and incurs a
unit production cost $c$ for \emph{every} offer, whether or not the buyer
accepts; successful sales replenish the budget.
This transforms seed selection into a \emph{sequential knapsack with network
externalities}: the policy must reason jointly about which buyers to skip,
in what order to approach them, and what price to offer, all under a live
cash-flow constraint.
We formalize this budget-constrained variant, introduce a
\emph{calibrated dynamic-programming} (Cal-DP) baseline that adapts classical
DP to replenishable-capital settings, and train a budget-aware learned policy
via imitation on Cal-DP trajectories followed by REINFORCE.

% ── P5: Contribution bullets C1–C4 ───────────────────────────────────────────
\smallskip
\noindent\textbf{Contributions.}
\begin{enumerate}[label=\textbf{C\arabic*.}, leftmargin=*, topsep=2pt,
                  itemsep=1pt, parsep=0pt]

  \item \textbf{Joint learning for revenue maximization.}
    We introduce the first end-to-end model for joint seed selection and
    personalized pricing in social networks.
    On Forest-Fire $n{=}1000$, the GNN+LSTM policy achieves
    $462.6$ %SOURCE: CLAUDE.md; comparison_20260623_131609.json
    vs.\ $460.0$ for the Babaei greedy baseline
    ($p{<}0.0001$, $20/20$ seeds). %SOURCE: CLAUDE.md §Session State

  \item \textbf{Zero-shot size and network generalization.}
    Trained on $n \le 440$, the model transfers without retraining to
    $n{=}2000$ ($+8.6\%$) %SOURCE: paper_gen_updated.json
    and to the real Rice-Facebook graph ($+40.5\%$), %SOURCE: paper_gen_updated.json
    demonstrating that the GNN+LSTM inductive bias learns
    topology-agnostic pricing strategies.

  \item \textbf{Budget-constrained formulation.}
    We formalize revenue maximization under replenishable finite capital as a
    sequential knapsack with network externalities, a practically
    relevant variant absent from prior work.

  \item \textbf{Calibrated DP + learned policy for the budget regime.}
    Our Cal-DP baseline achieves $99.8$ vs.\ $23.6$ for classical Greedy+Budget
    at $k{=}3$ ($4.2{\times}$); %SOURCE: dp_v3_full_curve_merged.json / dp_upgrade_eval.json
    the learned policy reaches $305.6$ vs.\ $7.4$ at $k{=}1$ ($41{\times}$). %SOURCE: unified_sweep.json / dp_upgrade_eval.json
    One honest finding: Cal-DP leads the learned policy at mid-$k$ on the
    dense Rice-Facebook graph — a result we report transparently as a
    regime-aware conclusion.

\end{enumerate}

% ── P6: Roadmap ───────────────────────────────────────────────────────────────
\smallskip
\noindent\textbf{Paper organization.}
\S\ref{sec:related} surveys related work.
\S\ref{sec:problem} formalizes unconstrained (\S\ref{sec:babaei}) and
budget-constrained (\S\ref{sec:budget}) revenue maximization.
\S\ref{sec:method} describes the GNN+LSTM model, training procedure, and
Cal-DP baseline.
\S\ref{sec:experiments} reports experimental results and the regime map.
\S\ref{sec:conclusion} concludes with limitations and future directions.
